\chapter{System Implementation} \label{chap:sysImplementation}

System Implementation explains how \textbf{RecruitSense} was engineered and deployed as a fully functional, multi-channel intelligent recruitment ecosystem, Applicant Tracking System (ATS), and Candidate Intelligence Platform. This chapter details the technical realization of the decoupled presentation tier (React Web and Flutter Mobile), the transactional backend services (Laravel 11), the relational database layer (MySQL 8.0), and the artificial intelligence and RAG microservice (FastAPI + ChromaDB).

\section{System Architecture Implementation}
RecruitSense is realized using the decoupled microservices architecture designed in Chapter \ref{chap:design} (Figure \ref{fig:sys_arch_detail}), cleanly isolating presentation, business orchestration, persistent storage, and artificial intelligence inference \cite{fielding2000rest, lewis2020rag, flutter2022google}.

\subsection{Architectural Overview}
The realized system comprises five coordinated layers:
\begin{itemize}
    \item \textbf{Web Presentation Layer:} A responsive Single Page Application (SPA) developed with React 18, Vite, and Tailwind CSS, delivering intuitive role-based web portals for Job Seekers, Corporate Recruiters, and Platform Administrators.
    \item \textbf{Mobile Presentation Layer:} A native-performing cross-platform mobile application engineered with Flutter and Dart, providing dedicated workspaces, candidate discovery, pipeline Kanban boards, and a live community networking feed on Android and iOS.
    \item \textbf{Application \& Gateway Layer:} A high-throughput RESTful API engine built with Laravel 11 (PHP 8.2+) that enforces business logic, role-based access control, hiring pipeline state transitions, domain validation, and secure authentication tokens.
    \item \textbf{Data Layer:} A normalized MySQL 8.0 relational database with structured 3NF schemas, foreign key cascade rules, and access-controlled disk storage for obfuscated resume documents.
    \item \textbf{AI Processing \& RAG Layer:} An asynchronous Python FastAPI microservice integrating PyPDF2 document stream parsing, regex skill tokenization, dense semantic vector embeddings, ChromaDB vector indexing, dynamic quiz generation, course recommendations, and local LLaMA-3.2 conversational RAG synthesis.
\end{itemize}

\subsection{Core Components}
\begin{itemize}
    \item \textbf{Web Frontend Component (React 18 SPA):} Structured around modular components using React 18 with Vite for rapid development and optimized bundling. Tailwind CSS provides responsive styling. Axios manages asynchronous HTTP communication with global interceptors.
    \item \textbf{Mobile Frontend Component (Flutter 3.x):} Developed using the Provider state management pattern. It handles reactive data binding, offline token caching, smooth screen transitions, and responsive mobile layouts.
    \item \textbf{Backend Component (Laravel 11 REST API):} Serves as the central orchestration authority. It validates incoming requests, verifies token identity via Laravel Sanctum, coordinates database transactions via Eloquent ORM, and communicates asynchronously with the AI microservice.
    \item \textbf{Data Component (MySQL 8.0 \& ChromaDB):} Persists relational domain records (users, companies, vacancies, applications, skill gaps) and high-dimensional resume chunk vector embeddings.
    \item \textbf{AI \& RAG Inference Component (FastAPI):} Ingests resumes, computes ATS match percentages, indexes dense vector embeddings into ChromaDB, executes semantic candidate searches, generates assessment quizzes, and synthesizes conversational RAG answers via LLaMA 3.2.
\end{itemize}

\subsection{Inter-Layer Communication}
Communication between the presentation clients (React Web and Flutter Mobile) and the Laravel backend is executed over stateless HTTPS RESTful API calls with JSON payloads. Authenticated operations pass cryptographic bearer tokens via the \texttt{Authorization: Bearer <token>} header. Communication between the Laravel backend and the Python FastAPI AI service occurs over internal HTTP POST endpoints (\texttt{http://127.0.0.1:8000/analyze-resume}, \texttt{http://127.0.0.1:8000/rag/ask}) with strict Pydantic payload validation and defensive timeout handling.

\section{Tools and Technologies Used}
The development of RecruitSense leveraged an enterprise-grade technology stack optimized for developer productivity, maintainability, and execution speed.

\subsection{Web Frontend Technologies}
\begin{itemize}
    \item \textbf{React 18:} Component-driven JavaScript library for reactive user interface rendering.
    \item \textbf{Vite 5:} High-performance build tool providing rapid Hot Module Replacement (HMR) and tree-shaken production bundles.
    \item \textbf{Tailwind CSS 3:} Utility-first CSS framework enabling consistent typography and dark/light theme management.
    \item \textbf{Axios:} Promise-based HTTP client with automated bearer token injection interceptors.
    \item \textbf{Lucide React:} Crisp vector icon set for intuitive visual affordances.
\end{itemize}

\subsection{Mobile Frontend Technologies}
\begin{itemize}
    \item \textbf{Flutter 3.x \& Dart:} Cross-platform framework compiling to ARM64 native machine code for Android and iOS.
    \item \textbf{Provider:} Reactive state management library providing clean separation of business logic and UI widgets.
    \item \textbf{Google Sign-In:} OAuth authentication integration for streamlined mobile user login.
    \item \textbf{Flutter Secure Storage / Shared Preferences:} Encrypted local key-value storage for bearer tokens and user preferences.
    \item \textbf{HTTP Package:} Asynchronous networking library configured with global JSON deserialization.
\end{itemize}

\subsection{Backend Technologies}
\begin{itemize}
    \item \textbf{Laravel 11 (PHP 8.2+):} Enterprise PHP framework providing an expressive MVC architecture, robust routing, and migration versioning.
    \item \textbf{Laravel Sanctum:} Secure authentication system managing cryptographically hashed Personal Access Tokens.
    \item \textbf{Eloquent ORM:} Active-record database abstraction layer ensuring prepared statements and complete parameterization.
    \item \textbf{Custom Middleware:} Modular request filters enforcing role authorization and corporate domain validation.
\end{itemize}

\subsection{AI, RAG, and Large Language Model Technologies}
\begin{itemize}
    \item \textbf{FastAPI (Python 3.11):} High-performance, asynchronous web framework built on Starlette and Pydantic with automated OpenAPI documentation \cite{ramirez2019fastapi}.
    \item \textbf{ChromaDB:} Open-source embedding vector database with persistent on-disk storage and HNSW indexing \cite{malkov2018hnsw}.
    \item \textbf{Dense Vector Embeddings (\texttt{nomic-embed-text} / Sentence Transformers):} Deep neural transformers mapping text passages into dense semantic vector representations \cite{reimers2019sentence}.
    \item \textbf{Ollama \& LLaMA 3.2 (3B):} High-throughput local Large Language Model runtime powering candidate RAG reasoning and dynamic quiz generation without cloud dependencies.
    \item \textbf{PyPDF2 \& python-docx:} Document stream decoders converting binary PDF and DOCX files into clean UTF-8 text strings.
\end{itemize}

\section{Application and Database Security Implementation}
RecruitSense enforces defense-in-depth security across all architectural layers to protect candidate data privacy, prevent unauthorized access, and ensure recruitment integrity.

\subsection{Authentication System}
User authentication is managed server-side. Upon registration or login, credentials are validated, passwords are encrypted using the Bcrypt hashing algorithm with a minimum cost factor of 12 rounds, and a cryptographically random 64-character Personal Access Token is generated. The plain-text token is returned to the client once, while its SHA-256 hash is persisted in the database. Subsequent requests authenticate via HTTP Bearer headers. Google OAuth is seamlessly verified and linked to user accounts.

\subsection{Authorization and Role Management (RBAC)}
Authorization is enforced through custom Laravel middleware that verifies user privileges before granting access to protected routes:
\begin{itemize}
    \item \textbf{Job Seeker Guard (\texttt{role:job\_seeker}):} Restricts access to candidate profile updates, resume uploads, application submissions, and personal saved jobs.
    \item \textbf{Corporate Recruiter Guard (\texttt{role:company}):} Restricts access to job vacancy publishing, candidate shortlisted rankings, pipeline stage updates, candidate RAG dialogs, and interview scheduling.
    \item \textbf{Platform Administrator Guard (\texttt{role:admin}):} Provides root access to corporate domain verification queues, content moderation tools, and global platform analytics.
\end{itemize}

\subsection{Database Security and Data Isolation}
Database security in RecruitSense is enforced through multiple complementary defensive layers:
\begin{itemize}
    \item \textbf{SQL Injection Prevention via Prepared Statements:} All database interactions are executed through Laravel's Eloquent ORM and PDO prepared statements with strict parameter binding, completely eliminating first-order and second-order SQL injection vulnerabilities.
    \item \textbf{Encrypted Credential and Key Storage:} User passwords are encrypted using Bcrypt with 12 hashing rounds. API secret keys and database connection credentials reside strictly within environment configuration files (\texttt{.env}) and are excluded from version control.
    \item \textbf{Multi-Tenant Data Isolation and Foreign Key Integrity:} Strict relational foreign key constraints with \texttt{ON DELETE CASCADE} or \texttt{RESTRICT} rules guarantee that recruiters can only access applicant data and job vacancies belonging strictly to their verified corporate entity.
    \item \textbf{Restricted Database User Privileges:} In production deployment, the database connection operates under a dedicated non-root MySQL user granted only minimum required Data Manipulation Language (DML) privileges (\texttt{SELECT}, \texttt{INSERT}, \texttt{UPDATE}, \texttt{DELETE}) without administrative schema modification rights.
    \item \textbf{Access-Controlled Resume Storage:} Candidate PDF resumes are stored outside the public document root directory with randomized UUID filenames, accessible only through authenticated server-side stream controllers that verify authorization before delivering file bytes.
\end{itemize}

\section{Processing Logic and Algorithms}

\subsection{Resume Parsing and Skill Extraction Pipeline}
The resume parsing pipeline converts unstructured PDF documents into structured candidate profiles:
\begin{enumerate}
    \item \textbf{Binary Ingestion \& PDF Stream Decoding:} PyPDF2 reads document byte streams page-by-page, stripping binary headers and invalid character encodings into raw UTF-8 text.
    \item \textbf{Text Sanitization \& Normalization:} The text stream undergoes lowercasing, whitespace collapsing, and non-printable symbol removal.
    \item \textbf{Boundary-Delimited Skill Tokenization:} Regular expressions with word boundaries (\texttt{\textbackslash b}) match terms against curated technical dictionaries (languages, frameworks, databases, cloud tools, DevOps):
    \begin{equation}
        \text{Pattern} = \text{\texttt{r'\textbackslash b(' + '|'.join(taxonomies) + ')\textbackslash b'}}
    \end{equation}
\end{enumerate}

\subsection{Dense Semantic Vector Embedding and ATS Composite Scoring}
The ATS scoring algorithm computes an explainable composite match percentage combining deterministic skill fulfillment and semantic vector similarity:
\begin{enumerate}
    \item \textbf{Skill Match Score ($S_{\text{skill}}$):}
    \begin{equation}
        S_{\text{skill}} = \frac{|\text{Candidate Skills} \cap \text{Required Job Skills}|}{|\text{Required Job Skills}|} \times 100
    \end{equation}
    \item \textbf{Dense Semantic Vector Similarity ($S_{\text{semantic}}$):}
    \begin{equation}
        S_{\text{semantic}} = \text{Cosine Similarity}(\vec{E}_{\text{resume}}, \vec{E}_{\text{job}}) \times 100
    \end{equation}
    where $\vec{E}$ represents the dense semantic embeddings generated by the embedding model.
    \item \textbf{Composite Weighted ATS Score ($S_{\text{ATS}}$):}
    \begin{equation}
        S_{\text{ATS}} = (0.60 \times S_{\text{skill}}) + (0.40 \times S_{\text{semantic}})
    \end{equation}
\end{enumerate}

\subsection{Retrieval-Augmented Generation (RAG) and LLaMA 3.2 Inference}
The Candidate RAG Chat engine enables recruiters to interrogate candidate resumes conversationally with guaranteed factual grounding:
\begin{enumerate}
    \item \textbf{Sliding Window Chunking:} Cleaned resume text is divided into overlapping segments of 150 words with a 30-word stride to preserve sentence context.
    \item \textbf{Dense Embedding Generation:} Each chunk is converted into a dense vector embedding using the embedding transformer.
    \item \textbf{ChromaDB Vector Indexing:} Vectors are indexed into persistent ChromaDB collections along with metadata (\texttt{resume\_id}, \texttt{chunk\_index}, \texttt{text}).
    \item \textbf{Semantic Query Retrieval:} When a recruiter submits a question $q$, ChromaDB retrieves the top-$k$ ($k=3$) nearest neighbor context chunks using cosine distance.
    \item \textbf{Evidence-Grounded LLaMA-3.2 Synthesis:} The retrieved chunks and user question are injected into an augmented prompt for the local LLaMA 3.2 model, synthesizing answers containing verbatim resume citations.
\end{enumerate}

\subsection{Skill Gap Course Recommendation Engine}
When a candidate application exhibits missing required skills ($\text{Missing} = \text{Required} \setminus \text{Matched}$), the system queries an educational knowledge base. For each missing competency, the course recommender generates structured course recommendations containing the course title, learning platform (Coursera, Udemy, edX, YouTube), estimated duration, and skill level (Beginner/Intermediate).

\subsection{Dynamic AI Assessment Quiz Generation}
To validate candidate competencies objectively, the FastAPI microservice integrates LLaMA 3.2 prompt engineering to synthesize technical multiple-choice assessment quizzes tailored specifically to vacancy requirements, generating structured JSON with 4 options, correct answer indices, and detailed explanations.

\section{GUI Design}
This section presents the realized graphical user interfaces for RecruitSense across the React Single Page Application (Web) and Flutter Mobile Application.

\subsection{Responsive Design Across Platforms}
RecruitSense is built with a mobile-responsive layout system using Tailwind CSS utility classes and Flutter reactive widget layouts, ensuring consistent visual affordance, adaptive typography, and high touch target accessibility across standard desktop monitors ($1920\times1080$), laptops ($1366\times768$), and mobile viewports.

\subsection{Registration and Login Interfaces}
User authentication interfaces provide seamless onboarding for Job Seekers and Corporate Recruiters. Figure \ref{fig:gui_reg_web} and Figure \ref{fig:gui_login_web} illustrate the desktop web registration and authentication screens, while Figure \ref{fig:gui_auth_mobile} presents the mobile authentication experience.

\begin{figure}[h!]
\centering
\fbox{\parbox[c][6.5cm][c]{0.88\linewidth}{\centering \textbf{Screenshot Placeholder: User Registration Page (Web)}\\\vspace{3mm}\small\textit{Reserved area for React Web User Registration interface with role selection and corporate domain validation}}}
\caption{User Registration Interface on React Web}
\label{fig:gui_reg_web}
\end{figure}

\begin{figure}[h!]
\centering
\fbox{\parbox[c][6.5cm][c]{0.88\linewidth}{\centering \textbf{Screenshot Placeholder: User Login Page (Web)}\\\vspace{3mm}\small\textit{Reserved area for React Web Authentication screen with credential and Google OAuth login}}}
\caption{User Login Interface on React Web}
\label{fig:gui_login_web}
\end{figure}

\begin{figure}[h!]
\centering
\fbox{\parbox[c][6.5cm][c]{0.88\linewidth}{\centering \textbf{Screenshot Placeholder: Mobile Registration and Login Screens (Flutter)}\\\vspace{3mm}\small\textit{Reserved area for Flutter Mobile Authentication and Splash Screens}}}
\caption{Mobile Registration and Login Interfaces (Flutter)}
\label{fig:gui_auth_mobile}
\end{figure}

\subsection{Job Discovery and Application Experience}
The job discovery portal enables candidates to search, filter, and inspect verified job postings. Figure \ref{fig:gui_browse_web} and Figure \ref{fig:gui_browse_mobile} show the web and mobile job browsing portals.

\begin{figure}[h!]
\centering
\fbox{\parbox[c][6.5cm][c]{0.88\linewidth}{\centering \textbf{Screenshot Placeholder: Public Job Discovery Portal (Web)}\\\vspace{3mm}\small\textit{Reserved area for React Web Job Search, Dynamic Filters, and Vacancy Card List View}}}
\caption{Public Job Discovery and Filtering Experience (Web)}
\label{fig:gui_browse_web}
\end{figure}

\begin{figure}[h!]
\centering
\fbox{\parbox[c][6.5cm][c]{0.88\linewidth}{\centering \textbf{Screenshot Placeholder: Mobile Job Discovery and Detail Views (Flutter)}\\\vspace{3mm}\small\textit{Reserved area for Flutter Mobile Job Search, Filter Drawers, and Job Detail Modals}}}
\caption{Mobile Job Discovery and Detail Views (Flutter)}
\label{fig:gui_browse_mobile}
\end{figure}

\subsection{Resume Upload, ATS Scoring, and Skill Gap Analysis}
When a candidate submits a PDF resume, the platform displays an interactive ATS match diagnostics modal. Figure \ref{fig:gui_ats_modal} illustrates the ATS score breakdown, matched skills, and personalized course recommendations.

\begin{figure}[h!]
\centering
\fbox{\parbox[c][6.5cm][c]{0.88\linewidth}{\centering \textbf{Screenshot Placeholder: ATS Diagnostics and Skill Gap Modal}\\\vspace{3mm}\small\textit{Reserved area for AI ATS Score Breakdown, Match Percentage, and Course Recommendations}}}
\caption{ATS Match Diagnostics and Skill Gap Analysis Modal}
\label{fig:gui_ats_modal}
\end{figure}

\subsection{Recruiter Dashboard, Candidate Shortlist, and Candidate RAG Chat}
The recruiter console provides corporate employers with candidate rankings sorted by computed ATS score, stage-by-stage pipeline Kanban progression, and conversational resume exploration via Candidate RAG. Figure \ref{fig:gui_recruiter_console} and Figure \ref{fig:gui_rag_chat} illustrate these interfaces.

\begin{figure}[h!]
\centering
\fbox{\parbox[c][6.5cm][c]{0.88\linewidth}{\centering \textbf{Screenshot Placeholder: Recruiter Candidate Shortlist and Kanban Board}\\\vspace{3mm}\small\textit{Reserved area for Candidate Ranking Table, ATS Match Badges, and Hiring Pipeline Kanban}}}
\caption{Recruiter Candidate Shortlist and Pipeline Management Interface}
\label{fig:gui_recruiter_console}
\end{figure}

\begin{figure}[h!]
\centering
\fbox{\parbox[c][6.5cm][c]{0.88\linewidth}{\centering \textbf{Screenshot Placeholder: Interactive Candidate RAG Chat Dialog}\\\vspace{3mm}\small\textit{Reserved area for Conversational Resume Q\&A Dialog with Semantic Text Evidence Highlights}}}
\caption{Candidate RAG Conversational Resume Intelligence Dialog}
\label{fig:gui_rag_chat}
\end{figure}

\subsection{Community Networking and Messaging Interfaces}
Figure \ref{fig:gui_messaging_feed} illustrates the mobile professional networking feed and real-time candidate-recruiter messaging interfaces.

\begin{figure}[h!]
\centering
\fbox{\parbox[c][6.5cm][c]{0.88\linewidth}{\centering \textbf{Screenshot Placeholder: Mobile Community Feed and In-App Messaging}\\\vspace{3mm}\small\textit{Reserved area for Flutter Mobile Networking Feed, Comments, and Real-Time Chat Screen}}}
\caption{Mobile Community Networking Feed and In-App Messaging Interface}
\label{fig:gui_messaging_feed}
\end{figure}

\subsection{Administrative Governance and Moderation Interface}
Figure \ref{fig:gui_admin_console} illustrates the platform administrator console for corporate domain approval, job posting moderation, and global platform telemetry.

\begin{figure}[h!]
\centering
\fbox{\parbox[c][6.5cm][c]{0.88\linewidth}{\centering \textbf{Screenshot Placeholder: Administrator Governance and Moderation Console}\\\vspace{3mm}\small\textit{Reserved area for Corporate Domain Verification Queue, Moderation Actions, and Audit Logs}}}
\caption{Platform Administrator Governance and Moderation Console}
\label{fig:gui_admin_console}
\end{figure}

\section{Deployment and Environment Implementation}
The implemented services are deployed across decoupled containerized runtime environments:
\begin{itemize}
    \item \textbf{Web Client Runtime:} The React SPA is bundled via Vite and served through an Nginx HTTP server configured for client-side routing.
    \item \textbf{Mobile Client Build:} The Flutter mobile codebase compiles to standalone Android APK/AAB packages and iOS application bundles.
    \item \textbf{Backend API Runtime:} The Laravel application runs on PHP-FPM 8.2 behind an Nginx reverse proxy with database connection pooling.
    \item \textbf{AI & RAG Service Runtime:} The FastAPI microservice runs via an asynchronous Uvicorn ASGI server managing local ChromaDB vector collections.
\end{itemize}

\section{Implementation Outcome}
The completed implementation of RecruitSense successfully fulfills all architectural specifications. The decoupled microservices tier delivers rapid, sub-second resume parsing, transparent ATS score formulations, evidence-grounded candidate Q\&A through ChromaDB vector retrieval, synchronized multi-platform access on web and mobile, and strict role governance across all user surfaces.

\section{Chapter Summary}
This chapter detailed the technical implementation of \textbf{RecruitSense}. It documented the multi-tier microservices architecture, detailed the technologies utilized across web, mobile, backend, and AI tiers, articulated application security safeguards, provided mathematical formulations for ATS scoring and RAG vector retrieval, elaborated the Flutter mobile application structure, and detailed the realized user interfaces across all operational surfaces. The verification and empirical evaluation of these implemented subsystems are presented in Chapter \ref{chap:testingEvaluation}.
